\documentclass[11pt]{article}
\usepackage[margin=1in]{geometry}
\usepackage{amsmath,amssymb,amsthm}
\usepackage[colorlinks=true,urlcolor=blue]{hyperref}
\newtheorem{theorem}{Theorem}
\newtheorem{proposition}[theorem]{Proposition}
\newtheorem{remark}[theorem]{Remark}
\newcommand{\one}{\boldsymbol 1}
\title{One unit edge edit preserves an optimal bisection;\\
two can change a unique optimum on half the vertices}
\author{Partitioning-limits session: exact elementary dichotomy}
\date{November 15, 2026}
\begin{document}
\maketitle
\begin{abstract}
For a graph with integer conductances, one unit insertion, deletion, or
conductance change always leaves a common old and new minimum bisection.
In particular, a unique old optimum remains optimal. The restriction is
sharp: two unit insertions can change the unique minimum bisection on
exactly half the vertices, even on connected simple unweighted graphs of
maximum degree five and arbitrarily large order. The same construction
works for two existing-edge doublings. These statements concern total
crossing weight under an exact cardinality constraint. They are elementary
structural limits, not a claim of a new spectral technique or historical
priority.
\end{abstract}

\section*{Definitions and the one-edit restriction}
All graphs share a fixed vertex set of even order $n$. A bisection has
exactly $n/2$ vertices on each side and is identified with its complement.
Its objective is the total crossing conductance. Conductance zero is
allowed when representing an absent edge. The feasible cut family does
not change during an edit. No assertion substitutes cardinality balance
for the degree-volume balance in normalized conductance.

\begin{theorem}[A common optimum under one unit edit]\label{thm:one}
Let $G$ have nonnegative integer conductances. Change one edge conductance
by a real $\delta$ with $|\delta|\le1$, leaving all conductances
nonnegative. For any fixed nonempty family $\mathcal F$ of cuts, the old
and new families of minimum-cost members have nonempty intersection.
If the old optimum is unique, it remains optimal after the change;
if $|\delta|<1$, it remains uniquely optimal.

In particular, a unit insertion or deletion in a simple unweighted graph,
or one existing unit edge doubled, cannot replace a unique minimum
bisection by a different unique minimum bisection.
\end{theorem}
\begin{proof}
For each feasible cut $S$, let $c(S)$ be its integer baseline cost and
let $\chi(S)\in\{0,1\}$ indicate whether it crosses the edited edge.
The new cost is $c(S)+\delta\chi(S)$.

If $\delta\ge0$, choose an old minimizer $S$ with smallest $\chi(S)$.
If $\chi(S)=0$, its cost is unchanged and no cut decreases in cost.
If $\chi(S)=1$, all old minimizers have this signature and their new
cost is $c(S)+\delta\le c(S)+1$. Every other cut starts at integer cost
at least $c(S)+1$ and does not decrease. Thus $S$ remains optimal.
For $\delta\le0$, choose an old minimizer with largest signature.
If its signature is one, it decreases by $|\delta|$, as much as any cut;
if its signature is zero, all old minima stay fixed and every other cut
can decrease by at most one from a baseline cost at least $c(S)+1$.
Again the chosen old optimum remains optimal.

When the old optimum $S$ is unique, every other feasible cut $T$ satisfies
$c(T)-c(S)\ge1$. Its relative updated cost is at least
$1-|\delta|$, proving persistence and the strict uniqueness assertion.
\end{proof}

\begin{remark}[What the integrality assumption supplies]
For an arbitrary weighted baseline with a unique optimum and objective
gap $\gamma>0$, the same comparison gives persistence for
$|\delta|\le\gamma$, and uniqueness for $|\delta|<\gamma$. Integer
conductances supply $\gamma\ge1$. Arbitrarily small noninteger gaps allow
weighted counterexamples with a single edit. The theorem does not assert
that every member of a tied old optimum family remains optimal.
\end{remark}

\section*{Sharpness with two edits}
\begin{theorem}[Two insertions give maximum disagreement]\label{thm:two}
For every odd $m\ge5$ there are connected simple unweighted graphs
$G_m,G'_m$ on $n=4m$ vertices, both of maximum degree at most five,
such that $G'_m$ is obtained from $G_m$ by inserting two edges. Both
graphs have a unique minimum bisection modulo complementation, and these
bisections satisfy
\[
 \min\{|S\triangle T|,|S\triangle T^c|\}=n/2.
\]
Equivalently, their normalized balanced sign vectors are orthogonal.
Instead of the two insertions, doubling two existing unit edges of the
same baseline graph has the identical conclusion.
\end{theorem}
\begin{proof}
Make four disjoint clouds $A,B,C,D$, each the square of the cycle $C_m$.
The edges of distance one form one Hamilton cycle, and the edges of
distance two form a second, edge-disjoint Hamilton cycle because $m$ is
odd. Thus the cloud is simple and four-regular. Every nontrivial vertex
subset cuts at least two edges in each Hamilton cycle, hence at least
four internal edges.

Add one $AB$ edge, one $CD$ edge, one $AC$ edge, and two $BD$ edges.
Use distinct external ports in each cloud. The three bisections that leave
every cloud whole have costs
\[
\begin{array}{c|ccc}
 &P=AB\mid CD&Q=AC\mid BD&R=AD\mid BC\\ \hline
G_m&3&2&5.
\end{array}
\]
There is a bisection of cost at most three. Any bisection splitting a
cloud cuts at least four internal edges and cannot be optimal. A
bisection leaving all equal clouds whole must place two clouds on each
side. Therefore the displayed list is exhaustive for optimizers, and
$Q$ is uniquely optimal modulo complement.

Now insert one additional $AB$ edge and one additional $CD$ edge, via
previously unused external ports. The three costs become
\[
\begin{array}{c|ccc}
 &P&Q&R\\ \hline
G'_m&3&4&7.
\end{array}
\]
The same internal four-edge argument rules out split clouds, so $P$ is
now the unique optimum. A cloud uses at most four external ports; these
exist since $m\ge5$. Every vertex has four internal edges and at most
one external edge. Both graphs are therefore simple, connected, and
of maximum degree at most five.

The sign patterns of $P$ and $Q$ on the four equal clouds are
$(+,+,-,-)$ and $(+,-,+,-)$. They have zero inner product and differ on
two clouds, establishing the asserted disagreement. Finally, doubling
the original $AB$ and $CD$ edges changes exactly the same three cut
costs. All other parts of the proof remain valid.
\end{proof}

\begin{proposition}[The one-unit endpoint is exact]
In the baseline graph of Theorem~\ref{thm:two}, increase only its
existing $AB$ edge by a real $t\ge0$. The unique minimum bisection is
$Q$ for $t<1$, is $P$ for $t>1$, and the optimum family at $t=1$ is
exactly $\{P,Q\}$. Thus one unit creates a tie, while every increase
larger than one selects the orthogonal unique bisection.
\end{proposition}
\begin{proof}
The three whole-cloud costs are $3,2+t,5+t$. An optimum has cost at
most three, so it cannot split a cloud, for any $t\ge0$.
\end{proof}

\section*{Relation to spectral stability}
These exact cut statements complement eigenvector perturbation results.
The first theorem supplies combinatorial persistence from integrality,
independently of spectral gaps. The second shows that two topological
edits already exhaust the largest possible disagreement of balanced
cuts. It does not assert that the Fiedler vectors of these particular
unweighted hosts recover the exact optima, or that their relative
Laplacian perturbation tends to zero. Those are separate statements.

The construction and proof are self-contained. The lack of a historical
priority claim is deliberate: the argument is an elementary consequence
of binary edge-crossing costs plus a four-cloud construction. Its value
here is the precise boundary it places on proposed one-unit
counterexamples and on interpretations of spectral instability.
\end{document}
